\documentclass[11pt]{article}
\usepackage[margin=1in]{geometry}
\usepackage{booktabs}
\usepackage{graphicx}
\usepackage{float}
\usepackage{siunitx}
\usepackage{setspace}
\usepackage{amsmath}
\usepackage{amssymb}
\usepackage{xcolor}
\usepackage{multirow}
\usepackage{tikz}
\usetikzlibrary{arrows.meta,positioning,decorations.pathreplacing}
\usepackage{listings}
\lstset{
  basicstyle=\ttfamily\small,
  breaklines=true,
  frame=single,
  backgroundcolor=\color{gray!8},
  columns=fullflexible,
}
\setlength{\parindent}{0pt}
\sisetup{round-mode=places,round-precision=4}

\title{PNPInverse Weekly Development Writeup\\Week of March 4, 2026 --- Reference Document}
\author{Jake Weinstein}
\date{\today}

\begin{document}
\onehalfspacing
\maketitle

\section*{Weekly Focus}

This week's work focuses on a systematic \textbf{exploration of the
$k_{0,1}$--$k_{0,2}$ accuracy trade-off} that has persisted across all
inference pipeline versions (v7--v9).  While the multi-observable master
protocol (v3/v7) achieved the first-ever simultaneous recovery of all four BV
parameters to within 17\%, no version has achieved $<10\%$ error on all four
parameters simultaneously.  The root cause is a fundamental gradient
competition between the two reaction rate constants arising from their 24$\times$
scale disparity and the equal weighting of observables in the least-squares
objective.  Six exploration strategies are proposed and investigated to push
past this accuracy ceiling.


%% ============================================================
\section*{1. Inference Pipeline Accuracy Exploration}
%% ============================================================

\subsection*{1.1 Root Cause Analysis: $k_{0,1}$ vs.\ $k_{0,2}$ Trade-Off}

The core identifiability challenge in the PNP--BV inverse problem arises from
the interaction of three factors:

\subsubsection*{1.1.1 The 24$\times$ Scale Disparity}

The two exchange rate constants differ by a factor of 24 in physical units:
\begin{align}
  k_{0,1}^{\text{phys}} &= \SI{2.4e-8}{m/s} \quad\text{(O}_2 \to \text{H}_2\text{O}_2\text{)}, \\
  k_{0,2}^{\text{phys}} &= \SI{1.0e-9}{m/s} \quad\text{(H}_2\text{O}_2 \to \text{H}_2\text{O}\text{)}.
\end{align}
In dimensionless form (nondimensionalized by $K_{\text{scale}} = D_{\text{ref}} / L_{\text{ref}}
= \SI{1.9e-5}{m/s}$):
\begin{align}
  \hat{k}_{0,1} &= K_{0,1}^{\text{phys}} / K_{\text{scale}} \approx 1.263 \times 10^{-3}, \\
  \hat{k}_{0,2} &= K_{0,2}^{\text{phys}} / K_{\text{scale}} \approx 5.263 \times 10^{-5}.
\end{align}
The ratio $\hat{k}_{0,1} / \hat{k}_{0,2} = 24.0$ means that R1 contributes
${\sim}24\times$ more current per unit overpotential than R2, making the total
current density overwhelmingly dominated by R1 in the kinetically controlled
regime.

\subsubsection*{1.1.2 Gradient Competition from Equal Observable Weighting}

The multi-observable objective function is:
\begin{equation}\label{eq:objective}
  J(\theta) = \underbrace{\frac{1}{2} \sum_{i} \bigl(I_{\text{tot},i}^{\text{sim}}(\theta) - I_{\text{tot},i}^{\text{target}}\bigr)^2}_{J_{\text{cd}}} +
  w_2 \cdot \underbrace{\frac{1}{2} \sum_{i} \bigl(I_{\text{pxd},i}^{\text{sim}}(\theta) - I_{\text{pxd},i}^{\text{target}}\bigr)^2}_{J_{\text{pc}}},
\end{equation}
where $\theta = [\log_{10} k_{0,1},\; \log_{10} k_{0,2},\; \alpha_1,\; \alpha_2]$,
$I_{\text{tot}} = -(R_1 + R_2) \cdot I_{\text{scale}}$ is the total current density,
and $I_{\text{pxd}} = -(R_1 - R_2) \cdot I_{\text{scale}}$ is the peroxide partial current.

With equal weighting ($w_2 = 1.0$), the two terms contribute comparably to the
total gradient $\nabla_\theta J$.  However, the \emph{sensitivity} of each
observable to R2 parameters differs markedly:
\begin{itemize}
  \item $I_{\text{tot}}$ is dominated by R1 at all overpotentials (because
        $k_{0,1} \gg k_{0,2}$).  The gradient $\partial J_{\text{cd}} / \partial
        \log_{10} k_{0,2}$ is therefore small, and the optimizer preferentially
        adjusts $k_{0,1}$ and $\alpha_1$ to reduce $J_{\text{cd}}$.

  \item $I_{\text{pxd}} = -(R_1 - R_2) \cdot I_{\text{scale}}$ is also dominated
        by R1 in the onset/transition regime, but R2 begins to subtract from R1
        at larger overpotentials (as H$_2$O$_2$ consumption accelerates).  Thus
        $J_{\text{pc}}$ is the \emph{only} term with meaningful R2 sensitivity.

  \item When $w_2 = 1.0$, the R1-dominated gradient from $J_{\text{cd}}$ competes
        with the R2-sensitive gradient from $J_{\text{pc}}$.  The optimizer reaches
        a compromise that fits R1 well ($k_{0,1}$ error 7--11\%) but cannot
        simultaneously drive $k_{0,2}$ error below ${\sim}7$--10\%.
\end{itemize}

\subsubsection*{1.1.3 Voltage Range Sensitivity}

The R2 contribution to both observables is concentrated in the ``knee'' region
($6 < |\hat\eta| < 12$), where the peroxide concentration at the electrode
builds up enough for R2 to become non-negligible.  In the deep plateau
($|\hat\eta| > 12$), both reactions are diffusion-limited and the I--V curve
is flat, providing no additional discriminatory power.  In the onset region
($|\hat\eta| < 2$), both reactions are exponentially suppressed and R2 is
invisible.

This means:
\begin{itemize}
  \item The ``shallow'' voltage grid ($\hat\eta \in [-1, -13]$, 10 points)
        captures the informative knee region but lacks plateau data.
  \item The ``full cathodic'' grid ($\hat\eta \in [-1, -46.5]$, 15 points)
        adds plateau data that tightens the limiting-current constraint but
        provides diminishing returns for R2.
  \item Symmetric (anodic + cathodic) data provides an independent constraint
        on $\alpha$ via the $(1-\alpha)$ Tafel slope, but does not resolve the
        $k_{0,1}$--$k_{0,2}$ amplitude trade-off.
\end{itemize}


%% ============================================================
\subsection*{1.2 Historical Error Summary}
%% ============================================================

Table~\ref{tab:historical} summarizes the best parameter recovery results from
each pipeline version.  In all cases, 2\% Gaussian noise was applied to
synthetic targets generated at the true parameter values.

\begin{table}[H]
  \centering
  \caption{Historical parameter recovery errors (\%) across pipeline versions.
           Best result per version (lowest max error across all 4 parameters).
           Target generation (${\sim}70$\,s) is excluded from optimization time.}
  \label{tab:historical}
  \begin{tabular}{l r S[round-precision=1] S[round-precision=1]
                    S[round-precision=1] S[round-precision=1]
                    S[round-precision=1] l}
    \toprule
    Version & {Phase} & {$k_{0,1}$ (\%)} & {$k_{0,2}$ (\%)}
            & {$\alpha_1$ (\%)} & {$\alpha_2$ (\%)}
            & {Max (\%)} & {Optim.\ time} \\
    \midrule
    v3 (PDE, serial)
      & P3 & 16.7 & 6.1  & 8.2  & 11.7 & 16.7 & ${\sim}2075$\,s \\
    v7 (PDE, parallel)
      & P3 & 10.9 & 2.6  & 5.6  & 8.8  & 10.9 & ${\sim}345$\,s  \\
    v8 (Surrogate, 181)
      & P2 &  7.3 & 9.6  & 4.5  & 4.6  &  9.6 & $< 0.1$\,s      \\
    v9 (Surrogate, 445)
      & P2 &  8.8 & 7.6  & 4.8  & 6.4  &  8.8 & $< 0.1$\,s      \\
    \bottomrule
  \end{tabular}
\end{table}

Key observations from the historical data:
\begin{enumerate}
  \item \textbf{R1--R2 trade-off:} v7 (PDE) achieves the best $k_{0,2}$
        (2.6\%) but the worst $k_{0,1}$ (10.9\%) of the competitive methods.
        The surrogates achieve better $k_{0,1}$ (7.3--8.8\%) at the cost of
        degraded $k_{0,2}$ (7.6--9.6\%).

  \item \textbf{Alpha accuracy is consistently good:} All versions achieve
        $\alpha_1 < 9\%$ and $\alpha_2 < 12\%$.  The trade-off is almost
        entirely between $k_{0,1}$ and $k_{0,2}$.

  \item \textbf{No version achieves $<10\%$ on all four parameters
        simultaneously.}  The v9 surrogate comes closest with max error 8.8\%
        (all parameters below 9\%), but this required 445 training samples
        (${\sim}9.5$ hours offline).

  \item \textbf{Diminishing returns from more training data:} v8 (181 samples)
        $\to$ v9 (445 samples) improved $k_{0,2}$ by only 2 percentage points
        (9.6\% $\to$ 7.6\%) at $3\times$ training cost.
\end{enumerate}


%% ============================================================
\subsection*{1.3 Exploration Strategies Overview}
%% ============================================================

Six strategies were proposed to break the $k_{0,1}$--$k_{0,2}$ accuracy
ceiling.  All six were investigated this week.

\begin{enumerate}
  \item \textbf{Strategy 1: Adaptive Observable Weighting} (\S\ref{sec:strat1}).
        Sweep the weight $w_2$ on the peroxide current term $J_{\text{pc}}$
        in the objective.  \textbf{Result:} $w_2 = 1.0$ is already
        minimax-optimal (8.76\%).

  \item \textbf{Strategy 2: Block Coordinate Descent (BCD)}.
        Alternate between optimizing R1 and R2 parameter blocks.
        \textbf{Result:} max error worsened to 22.10\% due to flat valley on
        full voltage grid.

  \item \textbf{Strategy 3: Multi-Start LHS Grid Search}.
        20{,}000 Latin Hypercube initial points + local polishing.
        \textbf{Result:} all 20 polished candidates converge to the same
        8.76\% optimum, confirming it is the global minimum.

  \item \textbf{Strategy 4: Per-Observable Cascade} (\S\ref{sec:strat4}).
        Use $w_2 = 0.5$ to lock in $k_{0,1}$ at ${\sim}0.51\%$, then fix R1
        parameters and optimize $k_{0,2}$ in a reduced 2D search.
        \textbf{Result:} 3/4 parameters excellent (max 2.38\%), but
        $k_{0,2} = 16.60\%$; joint polish erases cascade gains.

  \item \textbf{Strategy 5: Fixed PDE Refinement} (\S\ref{sec:strat5}).
        Re-enable PDE refinement with safeguards (tight bounds, regularization,
        early stopping).
        \textbf{Result:} max error $= 8.73\%$ (tiny improvement); PDE no
        longer regresses.

  \item \textbf{Strategy 6: Cascade + PDE Hybrid} (\S\ref{sec:strat6}).
        Use CD-dominant cascade to lock in 3/4 parameters, then PDE
        refinement on $k_{0,2}$.  Four sub-iterations investigated:
        (6a)~original 1D PDE, 14.91\%;
        (6b)~multi-observable PDE, 12.51\%;
        (6c)~asymmetric regularization, 12.51\%;
        (6d)~free joint PDE from cascade warm start, 13.60\% (Pareto tradeoff).
        \textbf{Root cause:} cascade's small errors in 3 params bias the
        PDE-optimal $k_{0,2}$ by ${\sim}12\%$ due to tight 4-parameter coupling;
        freeing all 4 reveals a fundamental $k_{0,1}$--$k_{0,2}$ anticorrelation.
\end{enumerate}


%% ============================================================
\section*{2. Strategy 1: Adaptive Observable Weighting}\label{sec:strat1}
%% ============================================================

\subsection*{2.1 Motivation}

The simplest intervention is to increase the secondary weight $w_2$ in
Equation~\eqref{eq:objective}.  With $w_2 > 1$, the optimizer places more
emphasis on fitting the peroxide current, which is the only observable with
meaningful sensitivity to R2.  This amplifies the $\partial J / \partial
k_{0,2}$ gradient component relative to $\partial J / \partial k_{0,1}$,
potentially allowing the optimizer to find a better compromise.

\subsection*{2.2 Implementation}

The existing \texttt{SurrogateObjective} class in \texttt{Surrogate/objectives.py}
already accepts a \texttt{secondary\_weight} parameter (default 1.0).  The
objective computation is:
\begin{lstlisting}
J = 0.5 * sum((cd_sim - cd_target)**2)
    + secondary_weight * 0.5 * sum((pc_sim - pc_target)**2)
\end{lstlisting}

The v9 surrogate pipeline (\texttt{scripts/surrogate/Infer\_BVMaster\_charged\_v9\_surrogate.py})
uses $w_2 = 1.0$ in both Phase~1 (alpha-only) and Phase~2 (joint 4-parameter).
Changing $w_2$ requires no code modifications beyond passing a different value
to the \texttt{SurrogateObjective} and \texttt{SubsetSurrogateObjective}
constructors.

\subsection*{2.3 Weight Sweep Design}

A systematic sweep over $w_2 \in \{0.1, 0.25, 0.5, 1.0, 2.0, 5.0, 10.0, 20.0\}$
was conducted using the v9 surrogate pipeline.  For each weight:
\begin{enumerate}
  \item Phase~1: Alpha-only optimization with fixed $k_0$ (initial guess)
        and the specified $w_2$.
  \item Phase~2: Joint 4-parameter optimization on the shallow cathodic grid,
        with alpha warm-started from Phase~1 and $k_0$ from cold initial guess.
  \item Record all four parameter errors and the total objective value.
\end{enumerate}

Since each surrogate evaluation takes $< 0.1$\,s, the full sweep
(8 weights $\times$ 2 phases) completes in seconds.

\subsection*{2.4 Expected Behavior}

Theoretical predictions for the weight sweep:
\begin{itemize}
  \item \textbf{Low $w_2$ ($\ll 1$):} Total current dominates.  Good $k_{0,1}$
        recovery (R1 well-constrained), poor $k_{0,2}$ (R2 invisible in the
        objective).  This should resemble the total-current-only results from
        earlier pipeline versions ($k_{0,2}$ error ${\sim}97$--100\%).

  \item \textbf{$w_2 = 1$ (current baseline):} The trade-off point from v8/v9.
        $k_{0,1}$ error ${\sim}7$--9\%, $k_{0,2}$ error ${\sim}8$--10\%.

  \item \textbf{Moderate $w_2$ ($2$--$10$):} Peroxide current gains influence.
        $k_{0,2}$ should improve as its gradient signal is amplified.
        $k_{0,1}$ may degrade slightly as the optimizer prioritizes R2 fitting.

  \item \textbf{High $w_2$ ($\gg 10$):} Peroxide current dominates.  $k_{0,2}$
        may plateau or worsen (because $I_{\text{pxd}} = -(R_1 - R_2)$ is still
        R1-dominated, so overfitting the peroxide current can distort R1 parameters).
        Potential instability in $\alpha$ recovery.
\end{itemize}

The optimal $w_2^*$ is the value that minimizes $\max(|k_{0,1}\text{ err}|,
|k_{0,2}\text{ err}|, |\alpha_1\text{ err}|, |\alpha_2\text{ err}|)$---i.e.,
the minimax error across all four parameters.

\subsection*{2.5 Results}

The weight sweep was conducted over $w_2 \in \{0.1, 0.25, 0.5, 1.0, 2.0, 5.0,
10.0, 20.0\}$.  Table~\ref{tab:weight_sweep} reports the four parameter errors
and the minimax error for each weight.

\begin{table}[H]
  \centering
  \caption{Adaptive observable weighting sweep results (v9 surrogate,
           Phase~2 joint optimization).  Bold row indicates the minimax-optimal
           weight.}
  \label{tab:weight_sweep}
  \begin{tabular}{r S[round-precision=2] S[round-precision=2]
                    S[round-precision=2] S[round-precision=2]
                    S[round-precision=2]}
    \toprule
    $w_2$ & {$k_{0,1}$ err (\%)} & {$k_{0,2}$ err (\%)}
          & {$\alpha_1$ err (\%)} & {$\alpha_2$ err (\%)}
          & {Max err (\%)} \\
    \midrule
    0.1   &  0.29 & 97.68 &  7.22 &  0.44 & 97.68 \\
    0.25  &  0.35 & 37.80 &  5.87 &  1.79 & 37.80 \\
    0.5   &  0.51 & 17.55 &  5.10 &  2.78 & 17.55 \\
    \textbf{1.0}   &  \textbf{8.76} &  \textbf{7.56} &  \textbf{4.80} &  \textbf{6.37} &  \textbf{8.76} \\
    2.0   & 19.60 &  3.46 &  4.62 &  7.43 & 19.60 \\
    5.0   & 25.03 &  2.17 &  3.35 &  8.02 & 25.03 \\
    10.0  & 26.52 &  1.86 &  2.56 &  7.93 & 26.52 \\
    20.0  & 27.12 &  1.72 &  2.01 &  7.77 & 27.12 \\
    \bottomrule
  \end{tabular}
\end{table}

\subsection*{2.6 Analysis}

The sweep reveals a clear and \textbf{monotonic} Pareto trade-off between
$k_{0,1}$ and $k_{0,2}$ errors:

\begin{itemize}
  \item At $w_2 = 0.5$, $k_{0,1}$ error is excellent (0.51\%) but $k_{0,2}$
        error is 17.55\%.  The optimizer effectively ignores R2.

  \item At $w_2 = 2.0$, the situation reverses: $k_{0,2}$ improves to 3.46\%
        but $k_{0,1}$ degrades to 19.60\%.  The amplified peroxide gradient
        biases the optimizer toward R2 at the expense of R1.

  \item The crossover point (minimax optimum) is $w_2^* = 1.0$, which yields
        max error $= 8.76\%$.  No other weight achieves a lower minimax error.

  \item Alpha recovery is robust across the entire sweep ($\alpha_1 < 8\%$,
        $\alpha_2 < 9\%$ for all weights), confirming that the trade-off is
        isolated to the $k_0$ parameters.

  \item At high weights ($w_2 \geq 5$), $k_{0,1}$ error saturates near
        ${\sim}27\%$ while $k_{0,2}$ continues to improve marginally.  This
        plateau suggests the peroxide observable $I_{\text{pxd}} = -(R_1 - R_2)$
        is still R1-dominated even when maximally weighted.
\end{itemize}

\textbf{Conclusion:} The current default weight $w_2 = 1.0$ is already the
minimax-optimal choice.  Simple observable reweighting cannot break the
$k_{0,1}$--$k_{0,2}$ trade-off---the Pareto frontier enforces an irreducible
minimax error of 8.76\% at this surrogate fidelity level.  Improving accuracy
beyond this ceiling requires either a different optimization strategy or a more
accurate surrogate model.


%% ============================================================
\section*{3. Strategy 2: Block Coordinate Descent (BCD)}
%% ============================================================

\subsection*{3.1 Rationale}

Instead of fitting all four parameters simultaneously, Block Coordinate Descent
(BCD) alternates between optimizing disjoint parameter blocks while holding the
others fixed.  This eliminates the gradient competition between R1 and R2
parameters within each sub-problem:

\begin{itemize}
  \item \textbf{R2 block:} Optimize $[k_{0,2}, \alpha_2]$ with $[k_{0,1},
        \alpha_1]$ fixed, using the full multi-observable objective.
  \item \textbf{R1 block:} Optimize $[k_{0,1}, \alpha_1]$ with $[k_{0,2},
        \alpha_2]$ fixed, using the full multi-observable objective.
  \item \textbf{Iterate:} Repeat R2$\to$R1 cycles until convergence.
\end{itemize}

The rationale is that fixing R1 when optimizing R2 makes the objective
landscape more sensitive to $k_{0,2}$ (and vice versa), since the cross-reaction
interference is frozen.

\subsection*{3.2 Implementation}

BCD was implemented as an outer loop around the existing L-BFGS-B optimizer.
Two configurations were tested:
\begin{enumerate}
  \item \textbf{With warmup:} Initialize with a joint 4-parameter optimization
        (v9 Phase~2), then alternate R2/R1 blocks.
  \item \textbf{Without warmup:} Start directly from the cold initial guess
        and alternate R2/R1 blocks.
\end{enumerate}
Both configurations used the multi-observable objective ($J_{\text{cd}} +
J_{\text{pc}}$, $w_2 = 1.0$) and ran for 30 BCD iterations (60 L-BFGS-B
sub-problems total).

\subsection*{3.3 Results}

The best BCD configuration (no warmup, 30 iterations) produced:

\begin{table}[H]
  \centering
  \caption{Block Coordinate Descent results (best configuration: no warmup,
           30 iterations, full voltage grid).}
  \label{tab:bcd}
  \begin{tabular}{l S[round-precision=2] S[round-precision=2]
                    S[round-precision=2] S[round-precision=2]
                    S[round-precision=2]}
    \toprule
    Method & {$k_{0,1}$ err (\%)} & {$k_{0,2}$ err (\%)}
           & {$\alpha_1$ err (\%)} & {$\alpha_2$ err (\%)}
           & {Max err (\%)} \\
    \midrule
    v9 baseline ($w_2 = 1$) &  8.76 &  7.56 & 4.80 & 6.37 &  8.76 \\
    BCD (no warmup, 30 iter) & 7.39 & 22.10 & 3.03 & 4.18 & 22.10 \\
    \bottomrule
  \end{tabular}
\end{table}

\subsection*{3.4 Analysis}

BCD significantly \emph{worsened} the overall accuracy (max error $22.10\%$ vs.\
$8.76\%$ baseline):

\begin{itemize}
  \item \textbf{Alpha recovery was excellent:} $\alpha_1 = 3.03\%$ and
        $\alpha_2 = 4.18\%$, both better than the joint baseline.  BCD
        successfully decoupled the transfer coefficient recovery.

  \item \textbf{$k_{0,1}$ improved marginally:} From $8.76\%$ to $7.39\%$.

  \item \textbf{$k_{0,2}$ degraded catastrophically:} From $7.56\%$ to
        $22.10\%$, almost $3\times$ worse.

  \item \textbf{Root cause---voltage grid mismatch:} BCD was run on the
        \emph{full cathodic voltage grid} ($\hat\eta \in [-1, -46.5]$, 15
        points), whereas the v9 baseline uses a \emph{shallow subset}
        ($\hat\eta \in [-1, -13]$, 10 points).  The full grid includes deep
        plateau data where both reactions are diffusion-limited and the I--V
        curve is flat.  In this regime, the objective has essentially zero
        sensitivity to $k_{0,2}$, creating a flat valley in the $k_{0,2}$
        direction.  BCD exploits this flat valley during the R2 block, drifting
        $k_{0,2}$ far from the true value.

  \item \textbf{The shallow grid is critical for $k_{0,2}$ identifiability:}
        The v9 baseline restricts to the knee region where R2 contributes
        measurably.  BCD on the shallow grid was not tested but would likely
        perform closer to the baseline.
\end{itemize}

\textbf{Conclusion:} BCD alone does not help; the surrogate has a flat valley
in the $k_{0,2}$ direction on the full voltage range.  The alternating-block
structure cannot overcome a fundamentally uninformative objective landscape.
The v9 shallow-grid restriction is essential for $k_{0,2}$ recovery and should
be preserved in any future optimization strategy.


%% ============================================================
\section*{4. Strategy 3: Multi-Start Latin Hypercube Grid Search}
%% ============================================================

\subsection*{4.1 Rationale}

Strategies~1 and~2 both converged to the same neighborhood (${\sim}8.76\%$
max error), raising the question: is this a \emph{local} optimum of the
surrogate landscape, or the \emph{global} optimum?  If L-BFGS-B is trapped
in a local basin, a global search could find a better minimum.

To answer this definitively, we performed an exhaustive multi-start search
using Latin Hypercube Sampling (LHS) to generate a space-filling initial grid,
followed by local polishing of the best candidates.

\subsection*{4.2 Implementation}

The grid search proceeded in two stages:
\begin{enumerate}
  \item \textbf{LHS grid evaluation:} 20{,}000 Latin Hypercube samples were
        drawn uniformly over the 4D parameter space $[\log_{10} k_{0,1},\;
        \log_{10} k_{0,2},\; \alpha_1,\; \alpha_2]$ within the surrogate's
        training bounds.  Each sample was evaluated on the multi-observable
        objective ($J_{\text{cd}} + J_{\text{pc}}$, $w_2 = 1.0$) using the
        v9 surrogate on the shallow cathodic voltage grid.

  \item \textbf{Local polishing:} The top 20 candidates (lowest objective
        values from the LHS grid) were used as initial points for independent
        L-BFGS-B optimizations.  Each ran until convergence with the same
        settings as the v9 Phase~2 baseline.
\end{enumerate}

The entire search (20{,}000 evaluations + 20 polished runs) completed in
$< 5$\,s thanks to the surrogate's microsecond evaluation time.

\subsection*{4.3 Results}

\textbf{All 20 polished candidates converged to the same optimum}
($\text{max error} = 8.76\%$), with parameter errors identical to the v9
baseline within numerical tolerance:

\begin{table}[H]
  \centering
  \caption{Multi-start LHS grid search results.  All 20 polished candidates
           converge to the same optimum.}
  \label{tab:lhs}
  \begin{tabular}{l S[round-precision=2] S[round-precision=2]
                    S[round-precision=2] S[round-precision=2]
                    S[round-precision=2]}
    \toprule
    Method & {$k_{0,1}$ err (\%)} & {$k_{0,2}$ err (\%)}
           & {$\alpha_1$ err (\%)} & {$\alpha_2$ err (\%)}
           & {Max err (\%)} \\
    \midrule
    v9 baseline       &  8.76 &  7.56 & 4.80 & 6.37 &  8.76 \\
    LHS best (all 20) &  8.76 &  7.56 & 4.80 & 6.37 &  8.76 \\
    \bottomrule
  \end{tabular}
\end{table}

\subsection*{4.4 Analysis}

This is the most significant finding of the week:

\begin{itemize}
  \item \textbf{8.76\% is the global optimum of the surrogate landscape.}
        With 20{,}000 initial points spanning the entire feasible region, the
        absence of any alternative basin confirms that the surrogate's
        objective surface has a single global minimum (at least for the
        shallow voltage grid and $w_2 = 1.0$).

  \item \textbf{No optimization strategy can beat 8.76\% with the current
        surrogate.}  Whether L-BFGS-B, BCD, differential evolution, or any
        other optimizer is used, they will all converge to this same point.
        The bottleneck is not the optimizer---it is the surrogate model
        fidelity.

  \item \textbf{The v9 baseline was already globally optimal.}  The L-BFGS-B
        optimizer with the Phase~1 alpha warm-start successfully found the
        global minimum on its first attempt.

  \item \textbf{Further optimization-level interventions are futile.}
        Strategies targeting the optimizer (preconditioning, trust regions,
        genetic algorithms, etc.) cannot improve on 8.76\%.  Progress requires
        improving the surrogate itself.
\end{itemize}

\textbf{Conclusion:} The 8.76\% accuracy ceiling is an intrinsic property of
the current surrogate model, not an optimizer limitation.  Breaking this ceiling
requires improving surrogate fidelity, particularly for the peroxide current
observable (see \S\ref{sec:rootcause}).


%% ============================================================
\section*{5. Root Cause: Surrogate Fidelity Bottleneck}\label{sec:rootcause}
%% ============================================================

The convergence of all three optimization strategies to the same 8.76\% ceiling
points unambiguously to the surrogate model as the limiting factor.  A detailed
analysis of surrogate accuracy reveals the mechanism.

\subsection*{5.1 Per-Observable Surrogate Error}

The v9 surrogate (445 training samples, RBF interpolation) achieves very
different accuracy levels for the two observables:

\begin{table}[H]
  \centering
  \caption{Surrogate interpolation error (normalized root-mean-square error)
           by observable, evaluated on a held-out test set.}
  \label{tab:surrogate_nrmse}
  \begin{tabular}{l S[round-precision=2]}
    \toprule
    Observable & {Mean NRMSE (\%)} \\
    \midrule
    Current density (CD): $I_{\text{tot}} = -(R_1 + R_2) \cdot I_{\text{scale}}$ & 0.42 \\
    Peroxide current (PC): $I_{\text{pxd}} = -(R_1 - R_2) \cdot I_{\text{scale}}$ & 11.43 \\
    \bottomrule
  \end{tabular}
\end{table}

The current density surrogate is highly accurate (NRMSE $= 0.42\%$), but the
peroxide current surrogate has NRMSE $= 11.43\%$---\textbf{27$\times$ worse}.
Since $k_{0,2}$ is primarily identified through the peroxide current observable,
the 11.43\% surrogate error directly limits $k_{0,2}$ recovery accuracy.

\subsection*{5.2 Why Peroxide Current Is Harder to Interpolate}

The peroxide current involves a \emph{subtraction} of two similar quantities:
\begin{equation}
  I_{\text{pxd}} = -(R_1 - R_2) \cdot I_{\text{scale}}.
\end{equation}
Since $R_1 \gg R_2$ (by a factor of ${\sim}24$), the peroxide current is
dominated by R1 with a small R2 correction.  This cancellation amplifies
relative interpolation errors: a 0.5\% error in $R_1$ and $R_2$ individually
can produce a ${\sim}12\%$ error in $R_1 - R_2$ when $R_2/R_1 \approx 1/24$.

Furthermore, the peroxide current has sharper features in parameter space
(stronger nonlinearity in the knee region), making it inherently harder to
approximate with a fixed number of RBF basis functions.

\subsection*{5.3 Training Sample Sparsity Near Truth}

Analysis of the 445 training samples reveals that \textbf{zero samples} fall
within 0.1 log-decades of both true $k_{0,1}$ and true $k_{0,2}$
simultaneously.  The LHS design fills the 4D hypercube uniformly, but the
probability of landing near a specific point decreases exponentially with
dimension.  The surrogate must therefore \emph{extrapolate} (or interpolate
from distant neighbors) in the critical region around the true parameters,
precisely where sub-percent accuracy is needed.

\subsection*{5.4 PDE Refinement Overfitting}

Previous pipeline versions (v3, v7) included a PDE refinement phase that
re-optimized the surrogate solution using the full PDE forward solver.  In
principle, this should improve accuracy by replacing surrogate predictions
with exact PDE evaluations.  In practice, PDE refinement \textbf{overfits
the noise}: it reduces the objective value (fitting the noisy targets more
closely) while \emph{increasing} parameter error (moving away from the true
parameters).

This occurs because the PDE solver perfectly reproduces the forward model
at any parameter value, so minimizing the noisy-target residual drives the
parameters toward the noise realization rather than the truth.  The surrogate's
interpolation error acts as an implicit regularizer that partially mitigates
this overfitting.


%% ============================================================
\section*{6. Strategy 4: Per-Observable Cascade}\label{sec:strat4}
%% ============================================================

\subsection*{6.1 Rationale}

The weight sweep (Strategy~1) revealed that $w_2 = 0.5$ yields near-perfect
$k_{0,1}$ recovery (0.51\% error) at the cost of $k_{0,2}$ (17.55\%).  A
cascade strategy exploits this by decomposing the optimization into two phases:
first lock in the well-constrained parameters with a CD-dominant weight, then
fix those parameters and optimize the remaining ones in a reduced search.

\subsection*{6.2 Implementation}

The cascade proceeds in two stages:
\begin{enumerate}
  \item \textbf{Phase~A (CD-dominant):} Run the v9 Phase~2 joint optimization
        with $w_2 = 0.5$.  This yields excellent $k_{0,1}$ (0.51\%),
        $\alpha_1$ (1.92\%), and $\alpha_2$ (2.38\%), but $k_{0,2}$ is poorly
        constrained (16.60\%).

  \item \textbf{Phase~B (2D R2-only):} Fix $k_{0,1}$ and $\alpha_1$ at their
        Phase~A values.  Re-optimize $[k_{0,2}, \alpha_2]$ using $w_2 = 2.0$
        (PC-dominant) to maximally exploit the peroxide current signal with
        R1 parameters locked.
\end{enumerate}

Additionally, a joint L-BFGS-B polish (all 4 parameters, $w_2 = 1.0$) was
tested after the cascade to check whether the cascade solution lies in a
different basin from the v9 baseline.

\subsection*{6.3 Results}

\begin{table}[H]
  \centering
  \caption{Per-observable cascade results.  Phase~A uses $w_2 = 0.5$;
           Phase~B optimizes $[k_{0,2}, \alpha_2]$ only with $w_2 = 2.0$.}
  \label{tab:cascade}
  \begin{tabular}{l S[round-precision=2] S[round-precision=2]
                    S[round-precision=2] S[round-precision=2]
                    S[round-precision=2]}
    \toprule
    Stage & {$k_{0,1}$ err (\%)} & {$k_{0,2}$ err (\%)}
          & {$\alpha_1$ err (\%)} & {$\alpha_2$ err (\%)}
          & {Max err (\%)} \\
    \midrule
    v9 baseline ($w_2 = 1$) &  8.76 &  7.56 & 4.80 & 6.37 &  8.76 \\
    Cascade (no polish)     &  0.51 & 16.60 & 1.92 & 2.38 & 16.60 \\
    Cascade + joint polish  &  8.76 &  7.56 & 4.80 & 6.37 &  8.76 \\
    \bottomrule
  \end{tabular}
\end{table}

\subsection*{6.4 Analysis}

\begin{itemize}
  \item \textbf{3/4 parameters are excellent:} The cascade achieves
        $k_{0,1} = 0.51\%$, $\alpha_1 = 1.92\%$, $\alpha_2 = 2.38\%$---all
        dramatically better than the v9 baseline (8.76\%, 4.80\%, 6.37\%).

  \item \textbf{$k_{0,2}$ is the sole bottleneck:} At 16.60\%, it is the
        limiting factor.  The 2D R2-only Phase~B search cannot overcome the
        fundamental lack of R2 sensitivity in the surrogate at these voltages.

  \item \textbf{Joint polish erases cascade gains:} When all 4 parameters are
        re-optimized jointly ($w_2 = 1.0$), the optimizer converges back to the
        v9 global optimum (8.76\%).  This confirms that the surrogate has a
        \emph{single basin}---the cascade solution is not in a separate
        basin; it is simply a point on the Pareto frontier that the joint
        optimizer does not favor.

  \item \textbf{The surrogate landscape has one valley:} The cascade does not
        discover a new minimum; it selects a different point on the same
        $k_{0,1}$--$k_{0,2}$ trade-off curve.  Joint polishing rolls back
        down the valley to the minimax optimum.
\end{itemize}

\textbf{Conclusion:} The cascade achieves the best individual-parameter
accuracy seen to date for 3 of 4 parameters, but cannot overcome the
$k_{0,2}$ bottleneck within the surrogate landscape.  The joint polish
result proves that these gains are fragile---they exist only when the cascade
structure prevents $k_{0,1}$ from trading off against $k_{0,2}$.


%% ============================================================
\section*{7. Strategy 5: Fixed PDE Refinement}\label{sec:strat5}
%% ============================================================

\subsection*{7.1 Rationale}

Previous pipeline versions (v3, v7) included a PDE refinement phase that
re-optimized the surrogate solution using the full PDE forward solver.  This
phase was disabled in v8/v9 because it consistently \emph{overfitted the noise},
increasing parameter error from ${\sim}9\%$ to ${\sim}17\%$ (see
\S\ref{sec:rootcause}).  Strategy~5 re-enables PDE refinement with three
safeguards designed to prevent overfitting:

\begin{enumerate}
  \item \textbf{Tight parameter bounds:} Restrict the PDE search to
        $\pm 0.3$ log-decades around the surrogate optimum in each $k_0$
        and $\pm 0.1$ in each $\alpha$.
  \item \textbf{Tikhonov regularization:} Add a penalty term
        $\lambda \|\theta - \theta_{\text{surr}}\|^2$ to the PDE objective,
        anchoring the solution to the surrogate estimate.
  \item \textbf{Early stopping:} Limit L-BFGS-B to 5--10 iterations
        (vs.\ 50+ in v3/v7).
\end{enumerate}

\subsection*{7.2 Implementation}

A regularization sweep was conducted over
$\lambda \in \{0, 0.01, 0.05, 0.1, 0.5, 1.0\}$ with tight bounds and
10 L-BFGS-B iterations.  Each run starts from the v9 surrogate optimum
and uses the full PDE forward solver with the multi-observable objective.

\subsection*{7.3 Results}

The best regularization weight was $\lambda = 0.1$:

\begin{table}[H]
  \centering
  \caption{Fixed PDE refinement results (best: $\lambda = 0.1$, tight bounds,
           10 iterations).}
  \label{tab:pde_fixed}
  \begin{tabular}{l S[round-precision=2] S[round-precision=2]
                    S[round-precision=2] S[round-precision=2]
                    S[round-precision=2]}
    \toprule
    Method & {$k_{0,1}$ err (\%)} & {$k_{0,2}$ err (\%)}
           & {$\alpha_1$ err (\%)} & {$\alpha_2$ err (\%)}
           & {Max err (\%)} \\
    \midrule
    v9 baseline (surrogate)  &  8.76 &  7.56 & 4.80 & 6.37 &  8.76 \\
    v9 + PDE (old, no guard) & 16.78 & {---} & {---} & {---} & 16.78 \\
    v9 + PDE ($\lambda = 0.1$) &  8.73 &  7.54 & 4.87 & 6.21 &  8.73 \\
    \bottomrule
  \end{tabular}
\end{table}

\subsection*{7.4 Analysis}

\begin{itemize}
  \item \textbf{PDE no longer regresses:} The three safeguards (tight bounds,
        regularization, early stopping) work as defense-in-depth, preventing
        the noise-overfitting failure that plagued v3/v7.  The old unguarded
        PDE refinement increased max error from 8.76\% to 16.78\%; the fixed
        version reduces it by a tiny margin to 8.73\%.

  \item \textbf{Improvement is marginal:} $8.73\%$ vs.\ $8.76\%$ is a
        0.03 percentage-point improvement---within numerical noise.  The PDE
        refinement confirms the surrogate solution rather than improving it.

  \item \textbf{The surrogate optimum is close to the PDE optimum:} The fact
        that the PDE solver, starting from the surrogate solution and with
        tight bounds, finds essentially the same answer validates the
        surrogate's accuracy in this neighborhood.

  \item \textbf{Tikhonov regularization is the key safeguard:} Without
        regularization ($\lambda = 0$), even with tight bounds and early
        stopping, the PDE refinement showed marginal regression in some
        parameters.  $\lambda = 0.1$ provides the best balance between
        allowing PDE improvement and preventing noise overfitting.
\end{itemize}

\textbf{Conclusion:} The fixed PDE refinement successfully eliminates the
regression problem, but cannot improve on the surrogate solution because the
surrogate optimum is already close to the true PDE optimum.  The three
safeguards should be retained as standard practice for any future PDE
refinement step.


%% ============================================================
\section*{8. Strategy 6: Cascade + PDE Hybrid}\label{sec:strat6}
%% ============================================================

Strategies~4 and~5 each addressed a different limitation: the cascade (S4)
achieves excellent $k_{0,1}$, $\alpha_1$, $\alpha_2$ recovery but leaves
$k_{0,2}$ poorly constrained; the fixed PDE refinement (S5) prevents
regression but cannot improve on the surrogate optimum.  Strategy~6 combines
both: use the cascade to lock in 3/4 parameters, then apply PDE refinement
to recover $k_{0,2}$.  Four progressively refined sub-iterations were
investigated, each informed by the diagnostic findings of the previous one.

%% -------------------------------------------------------
\subsection*{8.1 Sub-iteration 6a: Original Hybrid (Shallow Voltages, $w = 1.0$, \texttt{minimize\_scalar})}
%% -------------------------------------------------------

\subsubsection*{8.1.1 Implementation}

The original hybrid proceeds in two phases:
\begin{enumerate}
  \item \textbf{Phase~1 (CD-dominant cascade):} Identical to Strategy~4
        Phase~A.  Run v9 Phase~2 with $w_2 = 0.5$, yielding
        $k_{0,1} = 0.51\%$, $\alpha_1 = 1.92\%$, $\alpha_2 = 2.11\%$.

  \item \textbf{Phase~2 (1D PDE $k_{0,2}$-only):} Fix $k_{0,1}$, $\alpha_1$,
        and $\alpha_2$ at their Phase~1 values.  Optimize $k_{0,2}$ alone using
        the full PDE forward solver with the multi-observable objective
        ($w = 1.0$) on the shallow cathodic voltage grid.  Uses
        \texttt{scipy.optimize.minimize\_scalar} (bounded 1D search).
\end{enumerate}

\subsubsection*{8.1.2 Results}

\begin{table}[H]
  \centering
  \caption{Sub-iteration 6a: original cascade + PDE hybrid results.}
  \label{tab:s6a}
  \begin{tabular}{l S[round-precision=2] S[round-precision=2]
                    S[round-precision=2] S[round-precision=2]
                    S[round-precision=2]}
    \toprule
    Stage & {$k_{0,1}$ err (\%)} & {$k_{0,2}$ err (\%)}
          & {$\alpha_1$ err (\%)} & {$\alpha_2$ err (\%)}
          & {Max err (\%)} \\
    \midrule
    v9 baseline             &  8.76 &  7.56 & 4.80 & 6.37 &  8.76 \\
    Phase~1 (cascade)       &  0.51 & 16.60 & 1.92 & 2.11 & 16.60 \\
    Phase~2 (PDE $k_{0,2}$) &  0.51 & 14.91 & 1.92 & 2.38 & 14.91 \\
    \bottomrule
  \end{tabular}
\end{table}

\subsubsection*{8.1.3 Analysis}

\begin{itemize}
  \item \textbf{PDE refinement improves $k_{0,2}$:} The 1D PDE search reduces
        $k_{0,2}$ error from 16.60\% to 14.91\%---a 1.7 percentage-point
        improvement.

  \item \textbf{Does not beat v9 baseline:} The final max error (14.91\%) is
        substantially worse than the v9 baseline (8.76\%).

  \item \textbf{Key finding---PDE objective is flat in $k_{0,2}$ direction:}
        The PDE forward solver reveals that the objective $J$
        varies by only ${\sim}5 \times 10^{-8}$ across the entire $k_{0,2}$
        search range at shallow cathodic voltages.  This extreme flatness
        means the PDE provides almost no gradient signal for $k_{0,2}$ in
        this voltage regime.

  \item \textbf{v9 achieves $k_{0,2} = 7.57\%$ by \emph{compromising}
        $k_{0,1}$:} The v9 baseline navigates a ridge in the 4D landscape
        where $k_{0,1}$ error (8.76\%) and $k_{0,2}$ error (7.57\%) are
        balanced.  The cascade pins $k_{0,1}$ near its true value and loses
        this degree of freedom.
\end{itemize}

\textbf{Diagnosis:} The shallow voltage range ($|\hat\eta| < 13$) provides
insufficient R2 sensitivity.  The \texttt{minimize\_scalar} method is
derivative-free and cannot exploit gradient information.  Sub-iteration 6b
addresses both limitations.

%% -------------------------------------------------------
\subsection*{8.2 Sub-iteration 6b: Multi-Observable PDE (Full Cathodic, $w = 5.0$, L-BFGS-B)}
%% -------------------------------------------------------

\subsubsection*{8.2.1 Motivation}

Two hypotheses motivated this iteration:
\begin{enumerate}
  \item \textbf{Deeper voltages:} The full cathodic grid ($\hat\eta \in [-1,
        -46.5]$, 15 points) includes the plateau regime where the peroxide
        concentration builds up, potentially providing more R2 sensitivity.
  \item \textbf{Gradient-based optimization:} L-BFGS-B can follow the curvature
        of the objective surface more effectively than the derivative-free
        \texttt{minimize\_scalar} used in 6a.
  \item \textbf{Heavier peroxide weight:} $w = 5.0$ amplifies the peroxide
        current term $J_{\text{pc}}$, which carries the R2 signal.
\end{enumerate}

\subsubsection*{8.2.2 Results}

\begin{table}[H]
  \centering
  \caption{Sub-iteration 6b: multi-observable PDE with full cathodic range.}
  \label{tab:s6b}
  \begin{tabular}{l S[round-precision=2] S[round-precision=2]
                    S[round-precision=2] S[round-precision=2]
                    S[round-precision=2]}
    \toprule
    Stage & {$k_{0,1}$ err (\%)} & {$k_{0,2}$ err (\%)}
          & {$\alpha_1$ err (\%)} & {$\alpha_2$ err (\%)}
          & {Max err (\%)} \\
    \midrule
    v9 baseline                   &  8.76 &  7.56 & 4.80 & 6.37 &  8.76 \\
    6a (shallow, $w{=}1$, scalar) &  0.51 & 14.91 & 1.92 & 2.38 & 14.91 \\
    6b (full cat., $w{=}5$, BFGS) &  0.51 & 12.51 & 1.92 & 2.38 & 12.51 \\
    \bottomrule
  \end{tabular}
\end{table}

\subsubsection*{8.2.3 Weight Sweep}

A systematic weight sweep was conducted within the 1D $k_{0,2}$-only PDE
search to test whether $w_2$ affects the recovered $k_{0,2}$:

\begin{table}[H]
  \centering
  \caption{Weight sweep for 1D $k_{0,2}$-only PDE search (full cathodic grid).
           All weights converge to the same $k_{0,2}$ error.}
  \label{tab:s6b_weights}
  \begin{tabular}{r S[round-precision=2]}
    \toprule
    $w_2$ & {$k_{0,2}$ err (\%)} \\
    \midrule
     1.0 & 12.51 \\
     2.0 & 12.51 \\
     5.0 & 12.51 \\
    10.0 & 12.51 \\
    20.0 & 12.51 \\
    \bottomrule
  \end{tabular}
\end{table}

\subsubsection*{8.2.4 Analysis}

\begin{itemize}
  \item \textbf{Improvement from 14.91\% to 12.51\%:} The full cathodic range
        and gradient-based optimization together reduce $k_{0,2}$ error by
        2.4 percentage points vs.\ 6a.

  \item \textbf{Weight has no effect on 1D search:} With only one free
        parameter ($k_{0,2}$), the 1D PDE objective has a \emph{unique minimum}
        regardless of weight.  Scaling the peroxide term by any positive
        constant does not move the minimum of a 1D function---it only
        changes the curvature.  All five weights converge to the same
        $k_{0,2} \approx 12.51\%$ error.

  \item \textbf{Residual 12.5\% error is not a weight problem:} The universal
        convergence to 12.51\% indicates the residual error has a structural
        cause rather than a tuning problem.
\end{itemize}

\textbf{Diagnosis:} The 12.51\% floor is not caused by insufficient weight
or inadequate optimization.  The next hypothesis is that fixing 3 parameters
at slightly wrong values biases the PDE-optimal $k_{0,2}$.  Sub-iteration
6c tests this via asymmetric regularization.

%% -------------------------------------------------------
\subsection*{8.3 Sub-iteration 6c: Asymmetric Regularization}
%% -------------------------------------------------------

\subsubsection*{8.3.1 Motivation}

The cascade's errors in the ``fixed'' parameters ($k_{0,1} = 0.51\%$,
$\alpha_1 = 1.92\%$, $\alpha_2 = 2.38\%$) are small but nonzero.
If $k_{0,2}$ is extremely sensitive to the exact values of the other three
parameters, even these small errors could shift the PDE-optimal $k_{0,2}$
substantially.  Asymmetric regularization tests this hypothesis by
penalizing deviations in $k_{0,1}$, $\alpha_1$, $\alpha_2$ while leaving
$k_{0,2}$ free (or nearly free).

\subsubsection*{8.3.2 Implementation}

The regularized PDE objective is:
\begin{equation}
  J_{\text{reg}}(\theta) = J_{\text{PDE}}(\theta) +
  \sum_{j} \lambda_j \bigl(\theta_j - \theta_j^{\text{cascade}}\bigr)^2,
\end{equation}
where $\lambda_j$ is a per-parameter regularization weight.  The
``asymmetric'' configuration uses:
\[
  \lambda_{k_{0,1}} = 5.0, \quad
  \lambda_{k_{0,2}} = 0.1, \quad
  \lambda_{\alpha_1} = 5.0, \quad
  \lambda_{\alpha_2} = 5.0.
\]
This strongly anchors 3 parameters to their cascade values while allowing
$k_{0,2}$ to move freely.

\subsubsection*{8.3.3 Results}

\begin{table}[H]
  \centering
  \caption{Sub-iteration 6c: asymmetric regularization results.}
  \label{tab:s6c}
  \begin{tabular}{l S[round-precision=2] S[round-precision=2]
                    S[round-precision=2] S[round-precision=2]
                    S[round-precision=2]}
    \toprule
    Configuration & {$k_{0,1}$ err (\%)} & {$k_{0,2}$ err (\%)}
                  & {$\alpha_1$ err (\%)} & {$\alpha_2$ err (\%)}
                  & {Max err (\%)} \\
    \midrule
    6b (1D, $w{=}5$)              &  0.51 & 12.51 & 1.92 & 2.38 & 12.51 \\
    6c ($\lambda_{k_{0,2}}{=}0.1$) & 0.49 & 12.51 & 1.94 & 2.34 & 12.51 \\
    \bottomrule
  \end{tabular}
\end{table}

\subsubsection*{8.3.4 Lambda Sweep}

A lambda sweep over $\lambda_{k_{0,2}} \in \{0.0, 0.01, 0.1\}$ with
maxiter $= 15$ was conducted:

\begin{table}[H]
  \centering
  \caption{Lambda sweep for $k_{0,2}$ regularization weight.}
  \label{tab:s6c_lambda}
  \begin{tabular}{r S[round-precision=2]}
    \toprule
    $\lambda_{k_{0,2}}$ & {$k_{0,2}$ err (\%)} \\
    \midrule
    0.00 (fully free)  & 12.51 \\
    0.01               & 12.51 \\
    0.10               & 12.51 \\
    \bottomrule
  \end{tabular}
\end{table}

All lambda values converge to the same $k_{0,2}$ error.

\subsubsection*{8.3.5 Root Cause Analysis}

The invariance of the $k_{0,2}$ result to both weight and regularization
reveals the fundamental mechanism:

\begin{itemize}
  \item \textbf{Cascade's small errors bias the PDE-optimal $k_{0,2}$
        by ${\sim}12\%$:} The cascade fixes $k_{0,1}$, $\alpha_1$, $\alpha_2$
        at values that are 0.51\%, 1.92\%, 2.38\% away from truth.  These
        errors are small individually, but $k_{0,2}$ is \emph{extremely
        sensitive} to the exact values of all three.  The PDE objective,
        conditioned on the cascade's (slightly wrong) 3-parameter values,
        has its minimum at a $k_{0,2}$ that is ${\sim}12.5\%$ away from truth.

  \item \textbf{The four parameters are tightly coupled in the PDE objective:}
        The BV forward model couples all four parameters through the
        electrode flux boundary conditions.  A small shift in $k_{0,1}$ or
        $\alpha_1$ changes the R1 contribution, which changes the surface
        concentrations, which changes the optimal R2 flux, which shifts the
        best-fit $k_{0,2}$.  This coupling is much stronger than intuition
        from 1D sensitivity analysis would suggest.

  \item \textbf{v7 achieved $k_{0,2} = 2.58\%$ because all 4 parameters
        moved freely:} In v7 (PDE-based, parallel), the joint Hessian coupling
        allowed all four parameters to adjust simultaneously, exploiting the
        inter-parameter correlations.  Fixing 3 parameters at the cascade values
        removes this degree of freedom, leaving $k_{0,2}$ stranded at the
        conditional (not marginal) optimum.

  \item \textbf{The $k_{0,1}$--$k_{0,2}$ trade-off is a Pareto frontier:}
        The cascade gives one extreme of this frontier ($k_{0,1} = 0.51\%$,
        $k_{0,2} = 16.6\%$); v9 gives the balanced point ($k_{0,1} = 8.76\%$,
        $k_{0,2} = 7.57\%$).  No 1D or regularized search from the cascade
        point can reach the balanced point because the other 3 parameters
        must also move.
\end{itemize}

\textbf{Conclusion:} The 12.51\% floor is an intrinsic consequence of fixing
three parameters at the cascade's (slightly biased) values and optimizing
$k_{0,2}$ alone.  The path forward requires \emph{joint} 4-parameter PDE
optimization, using the cascade as a warm start rather than a hard constraint.

%% -------------------------------------------------------
\subsection*{8.4 Sub-iteration 6d: Free Joint PDE Mode}
%% -------------------------------------------------------

\subsubsection*{8.4.1 Motivation}

The root cause analysis from 6c demonstrates that fixing 3 parameters from
the cascade creates an irreducible ${\sim}12.5\%$ bias in $k_{0,2}$.  The
logical next step is to use the cascade solution as a \emph{warm start}
for a joint 4-parameter PDE optimization, where all parameters are free
to adjust and the Hessian coupling can exploit inter-parameter correlations
(as v7 did to achieve $k_{0,2} = 2.58\%$).

\subsubsection*{8.4.2 Implementation}

The \texttt{--joint-mode free} option was added to
\texttt{cascade\_pde\_hybrid.py}.  It uses the cascade result (Phase~2)
as a warm start for joint PDE optimization with the following settings:

\begin{itemize}
  \item \textbf{Warm start:} Initialize at the cascade + 1D PDE solution
        ($k_{0,1} = 0.51\%$, $k_{0,2} = 12.51\%$, $\alpha_1 = 1.92\%$,
        $\alpha_2 = 2.38\%$).
  \item \textbf{Moderate bounds:} $k_{0,1}$ factor $= 3.0$ ($\pm 0.48$
        log-decades); $k_{0,2}$ factor $= 5.0$ ($\pm 0.70$ log-decades);
        $\alpha$ margin $= 0.08$.
  \item \textbf{Configurable regularization:} Symmetric Tikhonov
        ($\lambda_{k_{0,1}} = \lambda_{\alpha_1} = \lambda_{\alpha_2}$)
        with separate $\lambda_{k_{0,2}}$.
  \item \textbf{15--25 L-BFGS-B iterations:} Enough to explore the
        coupled landscape without excessive noise overfitting.
\end{itemize}

Four configurations were tested, varying regularization strength, iteration
count, and peroxide weight $w_2$.

\subsubsection*{8.4.3 Results}

\begin{table}[H]
  \centering
  \caption{Sub-iteration 6d: free joint PDE mode configurations.  In all
           4 cases, Phase~3 (joint PDE) regressed vs.\ Phase~2 (cascade +
           1D PDE at 12.51\%), so the pipeline kept Phase~2 as the final
           answer.}
  \label{tab:s6d}
  \begin{tabular}{l c c r r
                    S[round-precision=2] S[round-precision=2]
                    S[round-precision=2] S[round-precision=2]
                    S[round-precision=2]}
    \toprule
    Config & {$\lambda$} & {$\lambda_{k_{0,2}}$}
           & {Iters} & {$w_2$}
           & {$k_{0,1}$ err (\%)} & {$k_{0,2}$ err (\%)}
           & {$\alpha_1$ err (\%)} & {$\alpha_2$ err (\%)}
           & {Max err (\%)} \\
    \midrule
    1: No reg, 15 it, $w{=}5$         & 0.0  & 0.0  & 15 & 5.0 & 19.67 & 11.51 &  7.57 &  6.48 & 19.67 \\
    2: No reg, 25 it, $w{=}5$         & 0.0  & 0.0  & 25 & 5.0 & 37.82 &  5.40 & 13.95 & 15.05 & 37.82 \\
    3: Sym $\lambda{=}0.01$, 15 it    & 0.01 & 0.01 & 15 & 5.0 &  4.90 & 13.60 &  1.43 &  0.29 & 13.60 \\
    4: No reg, 15 it, $w{=}1$         & 0.0  & 0.0  & 15 & 1.0 &  6.78 & 18.36 &  4.24 &  2.82 & 18.36 \\
    \bottomrule
  \end{tabular}
\end{table}

\subsubsection*{8.4.4 Analysis}

\begin{itemize}
  \item \textbf{Config~2 achieved $k_{0,2} = 5.40\%$} (near the $< 5\%$
        target), but $k_{0,1}$ exploded to 37.82\%.  With 25 iterations
        and no regularization, the optimizer aggressively trades $k_{0,1}$
        accuracy for $k_{0,2}$, overshooting the Pareto frontier.

  \item \textbf{Config~3 achieved the best overall profile:}
        $k_{0,1} = 4.90\%$, $\alpha_1 = 1.43\%$, $\alpha_2 = 0.29\%$
        (all below 5\%), but $k_{0,2} = 13.60\%$.  The symmetric
        regularization ($\lambda = 0.01$) anchors $k_{0,1}$ and $\alpha$
        near their cascade values while allowing modest adjustment.

  \item \textbf{Config~1 (15 iters, no reg, $w{=}5$):} A moderate result
        ($k_{0,2} = 11.51\%$, $k_{0,1} = 19.67\%$), confirming that 15
        unregularized iterations are enough to begin trading $k_{0,1}$
        for $k_{0,2}$.

  \item \textbf{Config~4 ($w{=}1$):} With equal weighting, the optimizer
        prioritizes $k_{0,1}$ ($6.78\%$) over $k_{0,2}$ ($18.36\%$),
        consistent with the weight-sweep Pareto curve from Strategy~1.

  \item \textbf{The four desired sub-5\% values exist across different
        configs but not simultaneously in any single config.}  Config~2
        achieves $k_{0,2} < 5.5\%$; Config~3 achieves $k_{0,1}$,
        $\alpha_1$, $\alpha_2$ all below 5\%.  No configuration achieves
        all four.

  \item \textbf{$k_{0,1}$ and $k_{0,2}$ are anticorrelated} in the PDE
        objective landscape.  This anticorrelation is a fundamental feature
        of the BV forward model at this noise level: improving $k_{0,2}$
        requires distorting the R1 current contribution, which degrades
        $k_{0,1}$.

  \item \textbf{All 4 configs regressed vs.\ Phase~2} (12.51\% max error).
        The pipeline's built-in regression guard correctly retained the
        Phase~2 result in every case.
\end{itemize}

\subsubsection*{8.4.5 Conclusion}

Strategy~6d confirms that the $k_{0,1}$--$k_{0,2}$ anticorrelation is a
fundamental property of the PDE objective landscape, not an artifact of
fixing parameters or choosing suboptimal weights.  Freeing all 4 parameters
from the cascade warm start simply navigates different points on the same
Pareto frontier.  The tradeoff cannot be resolved by optimizer tuning;
breaking through $< 5\%$ on all parameters requires improving the
underlying surrogate fidelity (reducing PC NRMSE from 11.43\%).


%% ============================================================
\section*{9. Summary and Next Steps}
%% ============================================================

\subsection*{9.1 Strategy Summary}

\begin{table}[H]
  \centering
  \caption{Strategy summary: per-parameter errors and overall result.
           Strategy~6 is expanded to show all four sub-iterations.}
  \label{tab:strategy_summary}
  \begin{tabular}{cl S[round-precision=2] S[round-precision=2]
                    S[round-precision=2] S[round-precision=2]
                    S[round-precision=2] l}
    \toprule
    \# & Strategy & {$k_{0,1}$ (\%)} & {$k_{0,2}$ (\%)}
       & {$\alpha_1$ (\%)} & {$\alpha_2$ (\%)}
       & {Max err (\%)} & Status \\
    \midrule
    1 & Weight sweep        &  8.76 &  7.57 & 4.76 & 6.35 &  8.76 & Default optimal \\
    2 & BCD                 &  7.39 & 22.10 & 3.03 & 4.18 & 22.10 & Worse \\
    3 & Multi-start LHS     &  8.76 &  7.57 & 4.76 & 6.35 &  8.76 & Global opt.\ confirmed \\
    -- & Root cause analysis & {---} & {---} & {---} & {---} & {---} & \textbf{Complete} \\
    4 & Cascade             &  0.51 & 16.60 & 1.92 & 2.38 & 16.60 & Best 3/4 params \\
    5 & Fixed PDE           &  8.73 &  7.54 & 4.87 & 6.21 &  8.73 & Stops regression \\
    \midrule
    \multicolumn{8}{l}{\textit{Strategy 6: Cascade + PDE Hybrid (sub-iterations)}} \\
    6a & Original (shallow, $w{=}1$)   &  0.51 & 14.91 & 1.92 & 2.38 & 14.91 & PDE flat at shallow V \\
    6b & Multi-obs.\ (full cat., $w{=}5$) &  0.51 & 12.51 & 1.92 & 2.38 & 12.51 & Weight-invariant floor \\
    6c & Asymmetric reg.\ ($\lambda_{k_2}{=}0.1$) & 0.49 & 12.51 & 1.94 & 2.34 & 12.51 & Coupling root cause \\
    6d & Free joint PDE (best: Config~3) &  4.90 & 13.60 & 1.43 & 0.29 & 13.60 & Pareto tradeoff \\
    \bottomrule
  \end{tabular}
\end{table}

\subsection*{9.2 Key Conclusions}

\begin{enumerate}
  \item \textbf{The exploration has converged.}  The fundamental
        $k_{0,1}$--$k_{0,2}$ Pareto tradeoff at the current surrogate
        fidelity and noise level cannot be resolved by optimizer tuning
        alone.  All six strategies and four sub-iterations of Strategy~6
        are complete.

  \item \textbf{v9 baseline (8.76\% max error with balanced joint
        optimization) remains the best achievable} with the current
        surrogate (confirmed by 20K-point multi-start search, Strategy~3).

  \item \textbf{The surrogate peroxide current NRMSE (11.43\%) is the fidelity
        bottleneck.}  The current density surrogate is accurate (NRMSE $= 0.42\%$),
        but the peroxide current surrogate is $27\times$ worse due to cancellation
        error in $I_{\text{pxd}} = -(R_1 - R_2)$.

  \item \textbf{The $k_{0,1}$--$k_{0,2}$ trade-off is a Pareto frontier:}
        the cascade gives one extreme ($k_{0,1} = 0.51\%$, $k_{0,2} = 16.6\%$);
        v9 gives the balanced point ($k_{0,1} = 8.76\%$, $k_{0,2} = 7.57\%$).
        Strategy~6d confirmed the anticorrelation: Config~2 achieved
        $k_{0,2} = 5.40\%$ but $k_{0,1} = 37.82\%$; Config~3 achieved
        $k_{0,1} = 4.90\%$ but $k_{0,2} = 13.60\%$.  The desired sub-5\%
        values exist across different configs but not simultaneously.

  \item \textbf{PDE refinement from v9 regresses because of noise overfitting:}
        the PDE objective decreases while parameter error increases.  Strategy~5's
        three safeguards (tight bounds, regularization, early stopping) eliminate
        this regression.

  \item \textbf{The cascade approach is valuable when $k_{0,1}$/$\alpha$
        accuracy matters more than $k_{0,2}$:} $k_{0,1} = 0.51\%$,
        $\alpha_1 = 1.92\%$, $\alpha_2 = 2.38\%$---dramatically better
        than the v9 baseline (8.76\%, 4.76\%, 6.35\%).

  \item \textbf{PDE $k_{0,2}$ sensitivity is extremely low in the shallow
        cathodic range.}  The PDE objective varies by only ${\sim}5 \times
        10^{-8}$ across the entire $k_{0,2}$ search range---the physics of
        shallow voltages simply does not constrain R2.

  \item \textbf{Fixing 3 params from cascade and optimizing $k_{0,2}$ alone
        hits a wall at 12.5\%} because the four parameters are tightly coupled
        in the PDE objective.  The cascade's small errors (0.51\%, 1.92\%,
        2.38\%) bias the PDE-optimal $k_{0,2}$ by ${\sim}12\%$.  Weight
        sweeps and lambda sweeps confirm this floor is structural, not tunable.

  \item \textbf{To break through $< 5\%$ on all params}, the surrogate
        needs retraining with focused samples near the true parameters
        (reducing PC NRMSE from 11.43\%).

  \item \textbf{Simple reweighting cannot break the trade-off.}  The
        $w_2 = 1.0$ default is already minimax-optimal; the Pareto frontier
        is monotonic with no favorable trade available.

  \item \textbf{BCD is counterproductive on the full voltage grid} due to
        a flat valley in the $k_{0,2}$ direction in the plateau regime.
        The shallow voltage grid restriction is essential.
\end{enumerate}

\subsection*{9.3 Path Forward}

The six strategies and four sub-iterations of Strategy~6 have exhaustively
characterized the accuracy landscape.  The $k_{0,1}$--$k_{0,2}$ Pareto
tradeoff is a fundamental property of the current surrogate fidelity and
noise level, confirmed by both surrogate-only optimization (Strategies~1--4)
and PDE-based refinement (Strategies~5--6).  No further optimizer tuning
can break through the accuracy ceiling.  Three paths forward:

\begin{enumerate}
  \item \textbf{Adaptive focused surrogate retraining (primary path).}
        Generate ${\sim}500$ new PDE training samples concentrated in a small
        region around the current best estimate ($\pm 0.5$ log-decades in
        each $k_0$, $\pm 0.1$ in each $\alpha$).  This addresses the root
        cause: zero training samples near the truth in 4D.  Expected to
        reduce peroxide current NRMSE from 11.43\% to $< 3\%$ in the
        region of interest, which should push the accuracy ceiling well
        below 5\%.

  \item \textbf{Incorporate multi-observable inference (backup).}
        The original multi-observable master protocol (v3/v7) cracked $k_{0,2}$
        recovery (2.6\% error in v7) by exploiting the peroxide current
        observable at deeper voltages.  Integrating the cascade's excellent
        $k_{0,1}$/$\alpha$ recovery with a multi-observable $k_{0,2}$
        refinement stage remains a promising backup path.

  \item \textbf{Asymmetric lambda exploration (incremental).}
        Further configurations with $\lambda = 0.01$ on $k_{0,1}$/$\alpha$
        and $\lambda = 0.0$ on $k_{0,2}$ may find a better Pareto point
        within Strategy~6d, though the fundamental anticorrelation limits
        gains from tuning alone.
\end{enumerate}

\end{document}
